\documentclass[tikz]{standalone}
\usepackage{avanti-color}
\usepackage{avanti-tikz4slides}
\usepackage{avanti-notations}

\tikzset{font=\large}
\tikzset{minimum-size/.style = {minimum height=0.8cm, minimum width=1.2cm}}
\tikzset{plain/.style = {draw=none, minimum-size}}
\tikzset{base/.style={smooth, thick, SolarizedBase01, draw=SolarizedBase01}}
\tikzset{hidden/.style = {rectangle, minimum-size, base}}
\tikzset{io/.style = {hidden, fill=SolarizedBase2}}
\tikzset{add/.style = {circle, minimum height=0.8cm, base, fill=SolarizedBase1}}

\tikzset{arrow/.style={->,>=stealth,base}}

\begin{document}

\begin{tikzpicture}

    \pgfmathsetmacro{\xinc}{2};
    \pgfmathsetmacro{\yinc}{1.5};

    \def \indexs {{1, 2, 3, 4, 5, "T"}};
    \def \names {{"\xv", "\av^{(1)}", "\av^{(2)}", "+", "\hat{y}"}};
    \def \styles {{"io", "hidden", "hidden", "add", "io"}};

    \foreach \i in {0, ..., 5}{

            \ifthenelse{\NOT \i = 4}{ % 非倒数第二列

                \pgfmathparse{\indexs[\i]};
                \edef \index {\pgfmathresult}; % 获取下标

                \foreach \j [count=\jj from -1] in {0, ..., 4}
                    {
                        \pgfmathparse{\styles[\j]};
                        \edef \style {\pgfmathresult}; % 获取样式

                        \ifthenelse{\j = 3}{ % 加号那一行不需要下标
                            \node [{\style}] (\i\j) at (\i*\xinc, \j*\yinc) {$\pgfmathparse{\names[\j]}\pgfmathresult$};
                        }{
                            \node [{\style}] (\i\j) at (\i*\xinc, \j*\yinc) {$\pgfmathparse{\names[\j]}\pgfmathresult_\index$};
                        }

                        \ifthenelse{\NOT \j = 0 \AND \NOT \j = 2}{
                            \draw [arrow] (\i\jj) -- (\i\j);
                        }{;}

                        \ifthenelse{\j = 2}{ % 画跳连边
                            \draw [arrow] (\i0) to [out=135, in=225] (\i2);
                        }{;}

                        \ifthenelse{\j = 3}{ % 画跳连边
                            \draw [arrow] (\i1) to [out=45, in=315] (\i3);
                        }{;}
                    }
            }{
                \foreach \j in {0, ..., 4}{
                        \node [plain,SolarizedBase01] (\i\j) at (\i*\xinc, \j*\yinc) {...};
                    }
            }
        }

    \foreach \i [count=\ii] in {0, ..., 5}{ % 画横向边
            \ifthenelse{\i < 5}{
                \draw [arrow] (\i1) -- (\ii1);
                \draw [arrow] (\ii2) -- (\i2);
            }{;}
        }

\end{tikzpicture}

\end{document}

